\documentclass[11pt,a4paper]{article}

\usepackage[T1]{fontenc}
\usepackage[utf8]{inputenc}
\usepackage[margin=2.5cm]{geometry}
\usepackage{amsmath}
\usepackage{amssymb}
\usepackage{booktabs}
\usepackage{graphicx}
\usepackage{listings}
\usepackage{xcolor}
\usepackage[numbers,sort&compress]{natbib}
\usepackage[hidelinks]{hyperref}

\newcommand{\always}{\ensuremath{\Box}}
\newcommand{\eventually}{\ensuremath{\Diamond}}
\newcommand{\code}[1]{\texttt{#1}}
\newcommand{\todo}[1]{\textcolor{red}{\textbf{[TODO: #1]}}}

\lstdefinelanguage{Promela}{
  morekeywords={active, proctype, init, chan, byte, bit, bool, short,
    int, mtype, pid, typedef, if, fi, do, od, break, goto, atomic,
    d_step, inline, ltl, never, assert, printf, skip, run, len, empty,
    nempty, full, nfull, timeout, else, of, select, for, in, unless,
    end, progress, accept},
  sensitive=true,
  morecomment=[l]{//}, morecomment=[s]{/*}{*/}, morestring=[b]"
}
\lstset{
  basicstyle=\ttfamily\footnotesize,
  keywordstyle=\bfseries,
  commentstyle=\itshape\color{gray},
  breaklines=true, showstringspaces=false,
  frame=single, framesep=5pt, xleftmargin=10pt, tabsize=2
}
\lstdefinestyle{promela}{language=Promela}
\lstdefinestyle{clang}{language=C}

\title{Model Checking Two Concurrent Programs with SPIN\\
       \large RW796 Software Verification and Analysis}
\author{Gabriella Scott \\ 25935453}
\date{\today}

\begin{document}
\maketitle

% Short. A paragraph or two at most: the two programs chosen, and the
% repository URLs, which the lecturer requires because the abstraction
% is marked with the implementation as reference.
\section*{Introduction}

\todo{One paragraph: which two programs were modelled and what the
analysis found.}

Implementation under analysis: \todo{repo URL}\\
Promela models, properties and verification records: \todo{repo URL}

% =====================================================================
\section{Program 1: FCFS Process Scheduler}

% ---------------------------------------------------------------------
\subsection{The Promela Model}

% Spec asks for: the design decisions, the relationship between model
% and implementation, the extent to which it is an appropriate
% abstraction and how that affects interpretation of results, and any
% major refinements made during the analysis.
%
% Cover, in prose, with short code excerpts only:
%
% 1. Configuration and why. NUM_PROCS 3, NUM_WORKERS 2, NUM_RESOURCES 2,
%    MAX_INSTR 3. Two workers is the smallest value at which any
%    interleaving exists; three processes is the smallest at which a non
%    trivial waiting queue exists; two resources allows a circular wait.
%    Ids start at 1 so 0 encodes NULL.
% 2. The mapping table below.
% 3. Locks as atomic. The omp critical plus q_locks[0] dequeue becomes
%    an atomic block \cite{holzmann2004}. Say explicitly where the
%    implementation holds no lock, since that is where the double
%    dispatch race lives, and note that those reads were deliberately
%    left outside any atomic block.
% 4. next_instruction as a bounded counter, with instruction type and
%    target resource chosen nondeterministically. Nondeterminism gives a
%    sound abstraction over all instruction sequences up to the bound
%    \cite{inggs_promela}.
% 5. Ghost arrays executing[] and in_ready[]. No counterpart in
%    manager.c, never read by a guard, so they do not change reachable
%    behaviour \cite{owicki1976}.
% 6. Omitted and why: dynamic arrivals and a_lock, the RR and priority
%    schedulers, names and strcmp, the nonexistent resource path,
%    logging, terminatedq as an explicit queue, detect_deadlock. For
%    each say what it costs. Two are not free: omitting arrivals removes
%    the terminate() race, and omitting detect_deadlock means deadlock
%    is checked by SPIN's invalid end state test rather than by the
%    program's own detector, which is a different question.
% 7. Refinements during the analysis: the blocking if with no
%    complementary guard, the else combined with i/o restriction,
%    active proctypes starting before init seeds ready_q, and the two
%    ghost arrays added when properties needed observables.
% 8. Limits on interpretation. The model is a hand built abstraction,
%    not a refinement, so a counterexample implies a real defect only
%    after checking by hand that the interleaving is permitted by the
%    implementation and by OpenMP \cite{openmp52}. The bounds mean a
%    property that holds establishes nothing outside them.

\todo{Model description.}

\begin{table}[htbp]
\centering
\caption{Mapping between the implementation and the model.}
\label{tab:mapping}
\begin{tabular}{p{4cm} p{4cm} p{6cm}}
\toprule
\textbf{\code{manager.c}} & \textbf{Model} & \textbf{Reason} \\
\midrule
Thread in \code{schedule\_fcfs} & \code{active [NUM\_WORKERS] proctype Worker} & \todo{} \\
\code{readyq}, \code{waitingq} & buffered channels & \todo{} \\
\code{terminatedq} & \code{pcb\_state[] == TERMINATED} & \todo{} \\
\code{pcb\_t.state} & \code{pcb\_state[]} & \todo{} \\
\code{next\_instruction} & \code{instr\_count[]} & \todo{} \\
\code{resource\_t.allocated} & \code{resource\_owner[]} & \todo{} \\
\code{q\_locks[0]}, \code{omp critical} & \code{atomic} block & \todo{} \\
\code{terminate()} & \code{check\_terminate} inline & \todo{} \\
\bottomrule
\end{tabular}
\end{table}

% Pair each excerpt with the manager.c fragment it abstracts, so the
% origin of the abstraction is visible. Three or four pairs is enough.

\begin{lstlisting}[style=clang,caption={Dequeue in \code{schedule\_fcfs}}]
#pragma omp critical
{
  omp_set_nest_lock(&q_locks[0]);
  running_p = dequeue_pcb(&readyq);
  omp_unset_nest_lock(&q_locks[0]);
}
\end{lstlisting}

\begin{lstlisting}[style=promela,caption={The same step in the model}]
\todo{paste the atomic dequeue block from model.pml}
\end{lstlisting}

% ---------------------------------------------------------------------
\subsection{Correctness Properties}

% For each property: informal statement, LTL formula, Promela form, and
% why it matters for the scheduler. Keep the justification tied to the
% implementation, not to the model.

\subsubsection*{\code{pcb\_mutex}: mutual exclusion over process control blocks}

\[ \always \bigwedge_{p=1}^{3} \mathit{executing}[p] \le 1 \]

\begin{lstlisting}[style=promela]
#define one_worker_per_pcb (executing[1] <= 1 && executing[2] <= 1 && executing[3] <= 1)
ltl pcb_mutex { [] one_worker_per_pcb }
\end{lstlisting}

% A PCB is the unit of scheduling. Two threads in the instruction loop
% for one PCB both advance next_instruction and both write state.
\todo{Justification.}

\subsubsection*{\code{no\_self\_wait}: no process waits for a resource it owns}

\[ \always \bigwedge_{p=1}^{3}
   \bigl( \mathit{waiting\_for}[p] = 0 \lor
          \mathit{resource\_owner}[\mathit{waiting\_for}[p]] \neq p \bigr) \]

% request_resource tests only resource->allocated == NULL and never
% compares allocated against the requester.
\todo{Promela form and justification.}

\subsubsection*{\code{waiting\_consistent}: the waiting state is consistent}

% pcb_state and waiting_for are two views of the same fact and a correct
% scheduler keeps them in agreement.
\todo{Formula, Promela form and justification.}

\subsubsection*{\code{no\_dup\_ready}: ready queue integrity}

% The only property expected to hold. Needs the in_ready ghost array,
% since the implementation keeps no count of queue membership.
\todo{Formula, Promela form and justification.}

\subsubsection*{\code{all\_terminate}: every process eventually terminates}

\[ \eventually \bigwedge_{p=1}^{3} \mathit{pcb\_state}[p] = \mathit{TERMINATED} \]

\todo{Justification.}

\subsubsection*{\code{ready\_dispatched}: a ready process is eventually dispatched}

\[ \always \bigl( \mathit{pcb\_state}[1] = \mathit{READY} \rightarrow
   \eventually\, \mathit{pcb\_state}[1] = \mathit{RUNNING} \bigr) \]

% Stated for process 1 only. Justify by symmetry: no process has a
% distinguishing attribute, so one instance is representative and the
% never claim stays small. Say that this is an argument, not a proof.
\todo{Justification.}

\subsubsection*{Deadlock}

% Not expressed in LTL. Checked with SPIN's invalid end state detection
% \cite{inggs_properties}, which asks whether the system halted where it
% should not have. That is a broader question than whether a circular
% wait exists \cite{coffman1971}.
\todo{Discussion.}

% ---------------------------------------------------------------------
\subsection{Verification Results and Analysis}

% The spec says not to merely reproduce SPIN output. One table, then
% interpretation.
%
% Setup, in two or three sentences: all properties are named ltl blocks
% in one model.pml selected with ./pan -a -N name, since pan silently
% uses the first claim if -N is omitted; an active never claim disables
% the invalid end state check, so the deadlock run uses a second binary
% compiled with -DNOCLAIM from the same pan.c \cite{spin_man}; MAX_INSTR
% is overridden on the command line rather than by editing the model.

\todo{Setup paragraph.}

\begin{table}[htbp]
\centering
\caption{Verification results, Spin 6.5.2. Depths and state counts as
reported by \code{pan}.}
\label{tab:results}
\begin{tabular}{l l l r r}
\toprule
\textbf{Property} & \textbf{Config} & \textbf{Result} & \textbf{Depth} & \textbf{States} \\
\midrule
\code{pcb\_mutex} & \code{MAX\_INSTR 3} & violated & 901  & 65\,957 \\
\code{no\_self\_wait} & \code{MAX\_INSTR 3} & violated & 379  & 188 \\
\code{waiting\_consistent} & \code{MAX\_INSTR 3} & violated & 962  & 69\,699 \\
\code{no\_dup\_ready} & \code{MAX\_INSTR 2} & holds    & n/a  & 20\,861\,451 \\
\code{all\_terminate} & \code{MAX\_INSTR 3}, \code{-f} & violated & 1069 & 493 \\
\code{ready\_dispatched} & \code{MAX\_INSTR 3}, \code{-f} & violated & 1088 & 11\,372\,991 \\
deadlock & \code{MAX\_INSTR 3}, \code{-DNOCLAIM} & violated & 504 & 67\,872 \\
\bottomrule
\end{tabular}
\end{table}

\subsubsection*{The double dispatch race}

% The headline finding, from pcb_mutex and waiting_consistent.
% - The interleaving in terms of manager.c: worker A dequeues p, blocks
%   it on a request, releases i_lock; before A reaches
%   if (running_p->state == WAITING) break;
%   worker B releases the resource, which moves p to readyq with state
%   READY; p is dequeued and run while A is still in its loop for p.
% - The cause: while (running_p->state == RUNNING) and the WAITING test
%   both read shared PCB state with no lock held.
% - Concrete harm: instr_count overshoots MAX_INSTR, which in the
%   implementation is next_instruction advanced twice for one
%   instruction, walking the pointer off the end of the list.
% - waiting_consistent shows the worse symptom: a PCB left WAITING, on
%   waiting_q, and owning a resource it will never release.
% - Fix: hold a lock across the state test and the branch, not only
%   across execution. Note the cost in lost parallelism.
% A figure of the interleaving would help here if you have space.

\todo{Analysis and fix.}

\subsubsection*{Self deadlock in \code{request\_resource}}

% Not a race, and the counterexample proves it: 379 deep, 188 states,
% one worker, executing[] never exceeds 1. Reachable single threaded.
% Cause: no ownership test, so a process re-requesting a resource it
% holds falls through with resource_found FALSE and is enqueued on
% waitingq to wait for itself. Note that the same branch also handles a
% resource that does not exist, which the model omits, so two different
% situations share one path. Fix: test ownership first and separate
% "not found" from "not available".

\todo{Analysis and fix.}

\subsubsection*{Liveness and fairness}

% Both liveness properties fail trivially without -f, at depth 14 with
% 25 states, in a cycle where init is never scheduled again so ready_q
% is never seeded. That is a modelling artefact, not a defect. Weak
% fairness excludes cycles in which a runnable process is passed over
% forever \cite{inggs_properties}, and costs a great deal: 25 states
% without -f against over eleven million with it for ready_dispatched.
% The fair all_terminate counterexample ends with process 1 holding
% resource 1 while waiting for resource 1, so it is the self wait defect
% resurfacing as a liveness failure rather than an independent finding.
% Say so, since a marker will notice.

\todo{Analysis.}

\subsubsection*{The invalid end state}

% Final state has instr_count = 4 with MAX_INSTR 3. A worker is blocked
% at the selection statement on instr_count because no guard is
% executable past the bound. In manager.c this is the instruction
% pointer walked off the end of the list by two threads. Argue whether
% modelling that as a blocked control point is faithful.

\todo{Analysis.}

\subsubsection*{The property that holds}

% What it establishes: within three processes, two workers, two
% resources and two instructions, no reachable state has a PCB twice in
% the ready queue, over 20,861,451 states with zero errors, and pan
% reported zero unreached statements in init and Worker, so the result
% is not an artefact of dead code \cite{inggs_spin}.
% What it does not establish: nothing outside those bounds, and nothing
% about the real linked list, since a channel cannot hold a duplicate
% reference the way a badly maintained list can. Be honest that the
% abstraction may have made this property easier to satisfy than it is
% in the implementation.
% Why MAX_INSTR 2: failing properties stop at the first counterexample,
% but a property that holds must exhaust the space, and the space grows
% by roughly 50 per instruction (420,705 at 1, 20,861,451 at 2, about
% 10^9 at 3, roughly 100 GB). The shortest possible duplicate needs one
% blocking request and one wakeup, so it fits within two instructions.
% That is an argument about the shortest witness, not a proof.

\todo{Analysis.}

% =====================================================================
\section{Program 2: \todo{name}}

\subsection{The Promela Model}
\todo{}

\subsection{Correctness Properties}
\todo{}

\subsection{Verification Results and Analysis}
\todo{}

% =====================================================================
\appendix
\section{Reproducing the verification results}

% The submission requires the exact commands and flags, and the ispin
% parameters if ispin was used. Record that ispin was not used.

\begin{lstlisting}[language=bash,caption={Build}]
spin -a model.pml
gcc -O2 -o pan pan.c
gcc -O2 -DNOCLAIM -o pan_noclaim pan.c
\end{lstlisting}

\begin{lstlisting}[language=bash,caption={Verification runs}]
./pan -a -N pcb_mutex
./pan -a -N no_self_wait
./pan -a -N waiting_consistent
./pan -a -N all_terminate    -f -m50000
./pan -a -N ready_dispatched -f -m50000 -w26
./pan_noclaim

spin -a -DMAX_INSTR=2 model.pml
gcc -O2 -o pan pan.c
./pan -a -N no_dup_ready -m20000 -w26
\end{lstlisting}

\todo{State the commit hash the results were produced from.}

\bibliographystyle{plainnat}
\bibliography{refs}

\end{document}